\documentclass[10pt,a4paper]{article}
\usepackage[margin=2.2cm]{geometry}
\usepackage{fontspec}
\setmonofont{DejaVu Sans Mono}[Scale=0.84]
\newfontfamily\unifont{DejaVu Serif}[Scale=0.9]
\usepackage{amsmath,amssymb,mathtools}
\setcounter{MaxMatrixCols}{20}
\usepackage{longtable,booktabs,array}
\usepackage{fancyvrb}
\usepackage[dvipsnames]{xcolor}
\usepackage[colorlinks=true,linkcolor=NavyBlue,urlcolor=NavyBlue]{hyperref}
\usepackage{microtype}
\usepackage{needspace}
\usepackage{tocloft}\setlength{\cftbeforesecskip}{0.45em}
\setlength{\parindent}{0pt}\setlength{\parskip}{0.55em}
\setlength{\emergencystretch}{8em}
\makeatletter
\renewcommand\section{\@startsection{section}{1}{\z@}{-3.5ex \@plus -1ex \@minus -.2ex}{2.3ex \@plus.2ex}{\normalfont\Large\bfseries\raggedright}}
\renewcommand\subsection{\@startsection{subsection}{2}{\z@}{-3.25ex\@plus -1ex \@minus -.2ex}{1.5ex \@plus .2ex}{\normalfont\large\bfseries\raggedright}}
\renewcommand\subsubsection{\@startsection{subsubsection}{3}{\z@}{-3.25ex\@plus -1ex \@minus -.2ex}{1.5ex \@plus .2ex}{\normalfont\normalsize\bfseries\raggedright}}
\makeatother
\renewcommand{\arraystretch}{1.15}
\allowdisplaybreaks
\begin{document}
\title{7+1 repository after the clean-up of 28 October 2026: provenance\\[0.4em]\large The task, the sources of every number, the decisions, the history, the building and the verification of the spreadsheet 7plus1\_status\_2026-10-28.xlsx and of the document on the state of the repository at commit 102f399}
\author{Repository: github.com/once-ere/7plus1 — prepared for Nash\\ Commit described: 102f399 of 28 October 2026, 13:15 PDT (branch main)\\ Document generated by tools/make\_docs.py on 2026-10-28}
\date{}
\maketitle
\begin{abstract}
This record states what was asked, where every number of the spreadsheet 7plus1\_status\_2026-10-28.xlsx comes from, every decision taken, how the version in the repository differs from the first version of the spreadsheet, how the spreadsheet and its three documents are built, and how they were checked: the 35 stored formula values against an evaluator of their own, the four commits, the 30 page counts and the changed documents against the repository, the 14 quoted formulas against their documents, and the counts of the twelve check programs and of the notebook against the outputs of this build. Every equation used is written here.
\end{abstract}
\tableofcontents
\Needspace{12\baselineskip}\section*{1. The task}\addcontentsline{toc}{section}{1. The task}
On 28 October 2026 the author asked for the following (summarized; the original wording is the author’s private working material and is not reproduced):

\begin{enumerate}
\item Remove the slightly loose lines that remained in the older PDF documents, then finish and push all stages.
\item Create a spreadsheet from the report on that work.
\item Confirm that all stages are finished, pushed and verified, or complete them.
\item Add the spreadsheet 7plus1\_status\_2026-10-28.xlsx to the repository and push it.
\end{enumerate}
The standing rules of the project apply: every documentation file has a companion provenance file and a symbol file, in Markdown and LaTeX, compiled to PDF; no file refers the reader to another documentation, provenance or symbol file; every equation that is created or used is written in the file that uses it. The repository’s .gitignore is respected.

\textbf{Interpretation.} In the repository the spreadsheet is a document like the others, so it gets a document that contains everything it contains (every sheet as a table, every formula with its value), a symbol dictionary and this provenance record. Its numbers are kept in a data file, \texttt{tools/\allowbreak{}status\_\allowbreak{}2026-10-28.json}, from which a program, \texttt{tools/\allowbreak{}make\_\allowbreak{}status\_\allowbreak{}xlsx.py}, writes the spreadsheet, so that it can be rebuilt and checked. The version in the repository is written for its readers: it describes the commit 102f399, refers to no conversation, and does not name the private files.

\Needspace{8\baselineskip}\subsection*{1.1 The signature: the author’s explanation of 27 October 2026}\addcontentsline{toc}{subsection}{1.1 The signature: the author’s explanation of 27 October 2026}
\textbf{Source.} On 27 October 2026 the author of the project explained that the xAct list $\{7,1,0\}$ declares seven space-like dimensions and one time-like dimension, while $\mathrm{diag}(+1,-1,\dots,-1)$ is simply the squares of the octonion basis, and asked that this be made clear in all documentation, examples and provenance records. \textbf{What it corrected.} The version of 7+1.nb pushed in the commits 52d8729 to dc2f178 of github.com/once-ere/7plus1 declared the xAct metric of section 12 with \texttt{DefMetric[\{1, 7, 0\}, …]}, reading the list as the signs of $\eta$; since then section 12 declares \texttt{DefMetric[\{7, 1, 0\}, g71[-a, -b], CD, …]}, exactly as the author’s notebook Einstein-Lovelock-7+1-Nash.nb does. \textbf{What was added.} Checks in 7+1.nb — section 4.7: $\eta$ is the squares of the octonion basis, $e_Ae_A=\eta_{AA}e_0$ and $\mathrm{Re}(e_Ae_B)=\eta_{AB}$; section 12.2: \texttt{SignatureOfMetric[g71]} $=\{7,1,0\}$ with \texttt{SignDetOfMetric[g71]} $=-1$, $\mathrm{Re}(e_Ae_B)=\eta_{AB}$, and $\mathrm{Re}(x\,x)=\eta_{AB}x^Ax^B$; section 17: the flat metric $\mathrm{diag}(+1,-1,-1,-1)$ of the world brane is the squares of the basis $e_0,e_1,e_2,e_3$ of the quaternion subalgebra — and the same statements in the Python programs \texttt{tools/\allowbreak{}ref71.py}, \texttt{tools/\allowbreak{}ref\_\allowbreak{}gravity.py} and \texttt{tools/\allowbreak{}ref\_\allowbreak{}sm.py}; the explanation itself in the text of 7+1.nb (sections 1.2, 1.4, 4.2, 4.7, 7, 11, 12.2, 13, 15.1, 17) and in every document, symbol file and provenance file of the repository. \textbf{What did not change.} No result. Compared output by output with the backup below, every output of 7+1.nb is the same, except the declared list itself, the new checks and the check counts (and, as in every run, the name of the private context and the measured times); in particular every component of every Lovelock tensor is unchanged. The components follow from the component matrix of the metric, and the sign of the determinant, \texttt{SignDetOfMetric}, is $-1$ for both lists. \textbf{Backup.} The evaluated notebook as it was before is kept in the workspace as \texttt{backups/\allowbreak{}7+1\_\allowbreak{}2026-10-27\_\allowbreak{}before\_\allowbreak{}signature.nb} (the folder backups/ is excluded from the repository by .gitignore), 1034401 bytes, SHA-256 \texttt{06094248\allowbreak{}8caf4f45\allowbreak{}120ffd37\allowbreak{}be0123e9\allowbreak{}6992a31c\allowbreak{}d7fdbde0\allowbreak{}9ee4dc4c\allowbreak{}ed4f5d3a}; \texttt{tools/\allowbreak{}verify\_\allowbreak{}7p1\_\allowbreak{}nb.py} compares every input, text and section cell of 7+1.nb with it and finds every cell present, each one unchanged or changed by this explanation, and it compares every output in the same way.

\textbf{The explanation, in full.} The 7+1 spacetime has one time-like dimension ($x_0$, the direction of the octonion unit $e_0=1$) and seven space-like dimensions ($x_1,x_2,x_3$ ordinary space, $x_4,\dots,x_7$ auxiliary space). Two different objects describe this, and neither may be replaced by the other. (i) The tensor package xAct, used in section 12 of 7+1.nb exactly as in the author’s notebook Einstein-Lovelock-7+1-Nash.nb, declares the metric with \texttt{DefMetric[\{7, 1, 0\}, g71[-a, -b], CD, …]}: the list $\{7,1,0\}$ declares seven space-like dimensions, one time-like dimension and no degenerate dimension — the numbers of entries $+1$ and $-1$ of $\eta_{71}=\mathrm{diag}(-1,+1,\dots,+1)$ (\texttt{η71}), the same spacetime written in the sign convention of 4+4.nb (space-like $=+1$, time-like $=-1$). (ii) The matrix $\eta=-\eta_{71}=\mathrm{diag}(+1,-1,\dots,-1)$ (\texttt{η17}) is simply the squares of the octonion basis: $e_0^2=+1$ for the unit $e_0=1$ (time) and $e_i^2=-1$ for the seven imaginary units $e_1,\dots,e_7$ (space); more generally it is the real part of the product of two basis units, and it gives the quadratic form of the octonions:
\begin{gather*}
\eta_{AB}=\mathrm{Re}(e_Ae_B),\qquad e_0^2=+1,\qquad e_i^2=-1\ \ (i=1,\dots,7),\\ \mathrm{Re}(x\,x)=x_0^2-x_1^2-\cdots-x_7^2=\eta_{AB}\,x^Ax^B\qquad (x=x^Ae_A\in\mathbb O).
\end{gather*}
$\eta$ gives the signs of the Clifford relation $\Gamma_A\Gamma_B+\Gamma_B\Gamma_A=2\,\eta_{AB}\,1_{16}$ of the real 16×16 Dirac matrices: time squares to $+1_{16}$ and each of the seven space directions to $-1_{16}$, exactly like the octonion units (in the notation $Cl(p,q)$, with $p$ generators squaring to $+1$ and $q$ to $-1$, the algebra $Cl(1,7)\cong\mathbb R(16)$). It also gives the signs of the metric components: $g_{00}>0$ along the time-like direction, negative squares along the space-like directions. Every component computed by 7+1.nb follows from the component matrix of the metric; the declared list fixes for xAct the sign of the determinant of the metric, \texttt{SignDetOfMetric[g71]} $=-1$, which is the same in both sign conventions, since $\det \eta=\det \eta_{71}=-1$. 7+1.nb checks \texttt{SignatureOfMetric[g71]} $=\{7,1,0\}$, \texttt{SignDetOfMetric[g71]} $=-1$, $\mathrm{Re}(e_Ae_B)=\eta_{AB}$ and $\mathrm{Re}(x\,x)=\eta_{AB}x^Ax^B$.

\Needspace{12\baselineskip}\section*{2. Where every number comes from}\addcontentsline{toc}{section}{2. Where every number comes from}
{\footnotesize
\begin{longtable}{>{\raggedright\arraybackslash}p{\dimexpr 0.2700\linewidth-2\tabcolsep\relax}>{\raggedright\arraybackslash}p{\dimexpr 0.4000\linewidth-2\tabcolsep\relax}>{\raggedright\arraybackslash}p{\dimexpr 0.3300\linewidth-2\tabcolsep\relax}}
\toprule
\textbf{data} & \textbf{source} & \textbf{how it was obtained} \\
\midrule\endfirsthead
\toprule
\textbf{data} & \textbf{source} & \textbf{how it was obtained} \\
\midrule\endhead
\endfoot
\bottomrule\endlastfoot
the commits: names, times, titles & the history of the repository & \texttt{git log}; the full names with \texttt{git rev-parse} \\
files changed, lines added, lines removed & the history of the repository & \texttt{git show --shortstat --format= <commit>} \\
what each commit contains & the commit messages & summarized \\
the documents, their kinds and what they cover & the list of documents of \texttt{tools/\allowbreak{}make\_\allowbreak{}docs.py} & — \\
pages & the PDFs of the commit 102f399 & \texttt{pdfinfo} \\
loose lines, first build & the logs of LaTeX of the first full build of the documents on 28 October 2026 (09:40 PDT, before commit 0654dcd) & the messages \texttt{Underfull \textbackslash{}\allowbreak{}hbox}, counted per document \\
loose lines after 0654dcd & the logs of the rebuild of a fresh clone of 0654dcd (10:10 PDT) & the same \\
loose lines, overflowing lines and errors now & the logs of the rebuild of a fresh clone of 102f399 (13:35 PDT) & the messages \texttt{Underfull \textbackslash{}\allowbreak{}hbox}, \texttt{Overfull \textbackslash{}\allowbreak{}hbox} and the lines that begin with “!” \\
changed by 93c500f, by 102f399 & the history of the repository & \texttt{git show --name-only --format= <commit>} \\
the rewritten paragraphs & the changes of the commits 0654dcd and 93c500f and the message of 93c500f & \texttt{git diff}, read paragraph by paragraph \\
the results of the twelve programs & their outputs in the rebuild of the fresh clone of 102f399 (build/docs/runs.json of that clone) & the exit codes, and the counts (2.1) below \\
the notebook & \texttt{tools/\allowbreak{}eval\_\allowbreak{}log\_\allowbreak{}7p1.txt} and the evaluated 7+1.nb & the record of the evaluation; the parse of the notebook by \texttt{tools/\allowbreak{}verify\_\allowbreak{}7p1\_\allowbreak{}nb.py} \\
the fresh clone & a clone of github.com/once-ere/7plus1 at 102f399, rebuilt from scratch & Markdown and LaTeX compared byte for byte, PDFs by \texttt{pdftotext} and \texttt{pdfinfo} \\
the private files & \texttt{git ls-files} of the fresh clone; \texttt{.gitignore}; \texttt{.git/\allowbreak{}info/\allowbreak{}exclude} on the author’s computer & — \\*
\end{longtable}}
The counts of the results of a program $P$, with $\ell$ running over the lines of its output and $\#$ counting the elements of a set:
\begin{align*}
N_{\rm True}(P)&=\#\{\ell:\ \ell\text{ ends in “: True” or in “(True)”}\},&&(2.1) \\
N_{\rm False}(P)&=\#\{\ell:\ \ell\text{ ends in “: False” or in “(False)”}\}
\end{align*}
(spaces between the colon and the word, and after the word, are allowed). A loose line is a line whose badness $b$ exceeds 1000, an overflowing line one that is wider than the line width $w$ by more than 0.1 pt even with all its spaces shrunk; with $t=w-w_0$ the stretch that a line of natural width $w_0$ needs, $s$ its total stretchability and $z$ its total shrinkability,
\begin{equation*}
b=\min\Big(10000,\ 100\Big(\frac ts\Big)^3\Big),\qquad \text{loose}\iff b>1000,\qquad \text{overflowing}\iff w_0-z-w>0.1\ \text{pt}.\qquad(2.2)
\end{equation*}
\Needspace{12\baselineskip}\section*{3. Decisions}\addcontentsline{toc}{section}{3. Decisions}
\begin{enumerate}
\item \textbf{A data file and a program}, not a spreadsheet made by hand: \texttt{tools/\allowbreak{}status\_\allowbreak{}2026-10-28.json} holds every number, and \texttt{tools/\allowbreak{}make\_\allowbreak{}status\_\allowbreak{}xlsx.py} writes the spreadsheet from it; the three documents read the same data file.
\item \textbf{Every total and every value of the sheet Summary is a formula} (35 formulas in all), so that the spreadsheet shows how each number follows from the tables.
\item \textbf{The values of the formulas are stored in the file.} openpyxl writes formulas without values; LibreOffice, run without a window, opens and saves the file, which computes and stores every value, so that a viewer that does not compute shows them too.
\item \textbf{The private files are not named}; they are described as the documents of the project describe them (a printout of an e-mail of the author on the radion; the author’s private notes on vielbeins in $n$ dimensions; a third private PDF).
\item \textbf{Times are PDT}, the time zone of the author’s computer, on which the commits were made.
\item \textbf{”Now” is the rebuild of a fresh clone of 102f399}, the most demanding test: only what is published is used.
\item \textbf{The documents reproduce the spreadsheet entirely} — every sheet as a table and every formula with its value — so that they can be read without it.
\item \textbf{Row 18 of the sheet Fixes} names the eight documents whose check counts changed in 102f399; a ninth, DIRAC16COMPLEX\_PROVENANCE, changed only in the SHA-256 of \texttt{tools/\allowbreak{}doc\_\allowbreak{}stage3\_\allowbreak{}common.py} that it records, and is marked changed in the sheet Documents.
\end{enumerate}
\textbf{The correction found while the data were gathered.} The function \texttt{parse\_\allowbreak{}ref} of \texttt{tools/\allowbreak{}doc\_\allowbreak{}stage3\_\allowbreak{}common.py} read a line \texttt{name : result} of the output of a program by splitting it at the FIRST “ : “. Two checks have names that contain “ : “ themselves — in \texttt{tools/\allowbreak{}ref\_\allowbreak{}coupled.py} a ratio \texttt{alpha2 E2 : alpha3 E3 : kappa eps -> 9/\allowbreak{}4 : -5/\allowbreak{}4 : 1}, in \texttt{tools/\allowbreak{}ref\_\allowbreak{}torus\_\allowbreak{}brane.py} the set \texttt{\#\{n in Z\textasciicircum{}4 : |n|\textasciicircum{}2 = k\}} — so their results were read as text, not as True, and the documents gave 49 and 37 checks instead of 50 and 38. Commit 102f399 splits a line that ends in “ : True” or “ : False” at the LAST “ : “ instead; the counts (2.1), which the spreadsheet uses, were right from the start.

\Needspace{12\baselineskip}\section*{4. The history of the spreadsheet}\addcontentsline{toc}{section}{4. The history of the spreadsheet}
A first version of the spreadsheet was made on 28 October 2026 from the same data and delivered to the author, who saved it in the working folder of the project (SHA-256 \texttt{179f155a\allowbreak{}c681ae0f\allowbreak{}3d7d283d\allowbreak{}c1e427e9\allowbreak{}dd4badbf\allowbreak{}2b3f2198\allowbreak{}090059f2\allowbreak{}a38d0b83}). The version in the repository was written by \texttt{tools/\allowbreak{}make\_\allowbreak{}status\_\allowbreak{}xlsx.py}. The two were compared cell by cell: every number and every computed value is the same; the differences are of wording and layout only:

\begin{itemize}
\item Summary: the subtitle names the commit described and its time; the first two items are called “Commit described (the latest when the spreadsheet was made)” and “Commits of 28 October 2026 up to the commit described” (before: “Latest commit”, “Commits since the radion stage began”).
\item Commits: three descriptions of contents are reworded, and the note says that the spreadsheet and its documents were added in the next commit, instead of referring to the conversation in which the first version was made.
\item Documents: the note defines the first build precisely (the first full build of 28 October 2026, before commit 0654dcd).
\item Fixes: the modifier letters of $\Theta^K$, $\sigma_r$ and $U^{\mathsf T}hU$ are written with \texttt{\textasciicircum{}} and \texttt{\_\allowbreak{}} (\texttt{Θ\textasciicircum{}K}, \texttt{σ\_\allowbreak{}r}, \texttt{U\textasciicircum{}T h U}), because some fonts lack them; row 18 names the eight documents and the ninth; the column Document is wider.
\item Checks: the note states the rule (2.1) by which the results are counted.
\item Private files: written in the third person and without the names of the files; two columns are wider.
\end{itemize}
\Needspace{12\baselineskip}\section*{5. How the spreadsheet and its documents are built}\addcontentsline{toc}{section}{5. How the spreadsheet and its documents are built}
\begin{enumerate}
\item \texttt{python3 tools/\allowbreak{}make\_\allowbreak{}status\_\allowbreak{}xlsx.py} reads the data file and writes the six sheets with openpyxl (3.1.5 when this record was generated, Python 3.13.17): Arial 10, headers in white on dark teal, the header row kept in view when scrolling (in Documents also the first column, in Summary the first four rows), landscape printing one page wide.
\item If LibreOffice is installed, the program converts the file with \texttt{soffice --headless --convert-to "xlsx:Calc MS Excel 2007 XML"} (with a private user profile in a temporary folder) and replaces it by the result: LibreOffice computes and stores the value of every formula and keeps the formulas. The spreadsheet in the repository records the application that saved it: \texttt{LibreOffice/\allowbreak{}24.2.7.2 Linux\_\allowbreak{}X86\_\allowbreak{}64 LibreOffice\_\allowbreak{}project/\allowbreak{}420 Build-2}.
\item The program reads the computed values back and fails if a formula has no value or an error value (a text that begins with “\#”).
\item \texttt{python3 tools/\allowbreak{}make\_\allowbreak{}docs.py} writes the document, its symbol dictionary and this record with \texttt{tools/\allowbreak{}doc\_\allowbreak{}status.py}, \texttt{tools/\allowbreak{}doc\_\allowbreak{}status\_\allowbreak{}symbols.py} and \texttt{tools/\allowbreak{}doc\_\allowbreak{}status\_\allowbreak{}provenance.py}, which take their data from \texttt{tools/\allowbreak{}doc\_\allowbreak{}status\_\allowbreak{}common.py}; that module reads the data file, rebuilds the workbook in memory and evaluates its formulas with an evaluator of its own (section 6.1).
\end{enumerate}
\Needspace{12\baselineskip}\section*{6. Verification}\addcontentsline{toc}{section}{6. Verification}
Every check of this section is recomputed whenever this record is generated.

\Needspace{8\baselineskip}\subsection*{6.1 The formulas of the spreadsheet}\addcontentsline{toc}{subsection}{6.1 The formulas of the spreadsheet}
The evaluator of \texttt{tools/\allowbreak{}doc\_\allowbreak{}status\_\allowbreak{}common.py} implements the functions and operators that the spreadsheet uses; with $R$ a range of cells, $v(c)$ the value of the cell $c$ and $x$, $y$ values:
\begin{align*}
\mathrm{SUM}(R)&=\sum_{c\in R,\ v(c)\ \text{a number}}v(c),\qquad \mathrm{COUNTA}(R)=\#\{c\in R:\ v(c)\ \text{not empty}\}, \\
\mathrm{COUNTIF}(R,x)&=\#\{c\in R:\ v(c)=x\}\quad(\text{texts compared without regard to case}), \\
x\,\&\,y&=\text{the text of }x\text{ followed by the text of }y .
\end{align*}
It evaluates the 35 formulas of the workbook that \texttt{tools/\allowbreak{}make\_\allowbreak{}status\_\allowbreak{}xlsx.py} makes from the data file; the values that LibreOffice stored in the spreadsheet equal them for \textbf{35 of 35} formulas.

{\footnotesize
\begin{longtable}{>{\raggedright\arraybackslash}p{\dimexpr 0.1300\linewidth-2\tabcolsep\relax}>{\raggedright\arraybackslash}p{\dimexpr 0.0700\linewidth-2\tabcolsep\relax}>{\raggedright\arraybackslash}p{\dimexpr 0.6000\linewidth-2\tabcolsep\relax}>{\raggedright\arraybackslash}p{\dimexpr 0.2000\linewidth-2\tabcolsep\relax}}
\toprule
\textbf{sheet} & \textbf{cell} & \textbf{formula} & \textbf{value} \\
\midrule\endfirsthead
\toprule
\textbf{sheet} & \textbf{cell} & \textbf{formula} & \textbf{value} \\
\midrule\endhead
\endfoot
\bottomrule\endlastfoot
Summary & B5 & \texttt{=Commits!A5} & 102f399 \\
Summary & B6 & \texttt{=COUNTA(Commits!A2:A5)} & 4 \\
Summary & B7 & \texttt{=Commits!E6} & 206 \\
Summary & B8 & \texttt{=COUNTA(Documents!A2:A31)} & 30 \\
Summary & B9 & \texttt{=Documents!D32} & 478 \\
Summary & B10 & \texttt{=Documents!E32} & 40 \\
Summary & B11 & \texttt{=Documents!F32} & 12 \\
Summary & B12 & \texttt{=Documents!G32} & 0 \\
Summary & B13 & \texttt{=Documents!H32} & 0 \\
Summary & B14 & \texttt{=Documents!I32} & 0 \\
Summary & B15 & \texttt{=COUNTIF(Fixes!D2:D19,"Display")+COUNTIF(Fixes!D2:D19,"Code block")+COUNTIF(Fixes!D2:D19,"Reworded")} & 16 \\
Summary & B16 & \texttt{=COUNTA(Checks!A2:A13)} & 12 \\
Summary & B17 & \texttt{=COUNTIF(Checks!C2:C13,0)} & 12 \\
Summary & B18 & \texttt{=Checks!D14} & 504 \\
Summary & B19 & \texttt{=Checks!E14} & 0 \\
Summary & B20 & \texttt{=Checks!B17} & 479 \\
Summary & B21 & \texttt{=Checks!B18} & 479 \\
Summary & B22 & \texttt{=Checks!B21} & 0 \\
Summary & B23 & \texttt{=Checks!B30\&" of "\&Checks!B29} & 60 of 60 \\
Summary & B24 & \texttt{=Checks!B32\&" of "\&Checks!B31} & 30 of 30 \\
Summary & B25 & \texttt{=COUNTIF('Private files'!C2:C5,"Yes")} & 0 \\
Commits & E6 & \texttt{=SUM(E2:E5)} & 206 \\
Commits & F6 & \texttt{=SUM(F2:F5)} & 25,646 \\
Commits & G6 & \texttt{=SUM(G2:G5)} & 5,357 \\
Documents & D32 & \texttt{=SUM(D2:D31)} & 478 \\
Documents & E32 & \texttt{=SUM(E2:E31)} & 40 \\
Documents & F32 & \texttt{=SUM(F2:F31)} & 12 \\
Documents & G32 & \texttt{=SUM(G2:G31)} & 0 \\
Documents & H32 & \texttt{=SUM(H2:H31)} & 0 \\
Documents & I32 & \texttt{=SUM(I2:I31)} & 0 \\
Documents & J32 & \texttt{=COUNTIF(J2:J31,"Yes")} & 11 \\
Documents & K32 & \texttt{=COUNTIF(K2:K31,"Yes")} & 9 \\
Checks & B14 & \texttt{=COUNTIF(C2:C13,0)\&" of 12 programs exit with code 0"} & 12 of 12 programs exit with code 0 \\
Checks & D14 & \texttt{=SUM(D2:D13)} & 504 \\
Checks & E14 & \texttt{=SUM(E2:E13)} & 0 \\*
\end{longtable}}
\Needspace{8\baselineskip}\subsection*{6.2 The data file against the repository}\addcontentsline{toc}{subsection}{6.2 The data file against the repository}
{\footnotesize
\begin{longtable}{>{\raggedright\arraybackslash}p{\dimexpr 0.5500\linewidth-2\tabcolsep\relax}>{\raggedright\arraybackslash}p{\dimexpr 0.4500\linewidth-2\tabcolsep\relax}}
\toprule
\textbf{check} & \textbf{result} \\
\midrule\endfirsthead
\toprule
\textbf{check} & \textbf{result} \\
\midrule\endhead
\endfoot
\bottomrule\endlastfoot
the full names of the four commits & 4 of 4 agree \\
files changed, lines added, lines removed (\texttt{git show --shortstat}) & 4 of 4 agree: 0654dcd 122, 24996, 4890; 54862f0 12, 66, 19; 93c500f 44, 210, 93; 102f399 28, 374, 355 \\
the page counts of the 30 PDFs at 102f399 (\texttt{git show 102f399:<document>.pdf}, \texttt{pdfinfo}) & 30 of 30 agree; 478 pages in all \\
the documents changed by 93c500f (\texttt{git show --name-only}) & 11 marked, 11 changed: the same documents \\
the documents changed by 102f399 (\texttt{git show --name-only}) & 9 marked, 9 changed: the same documents \\*
\end{longtable}}
\Needspace{8\baselineskip}\subsection*{6.3 The quoted formulas against their documents}\addcontentsline{toc}{subsection}{6.3 The quoted formulas against their documents}
Each formula of the rewritten paragraphs (section 6.5) must occur, character for character, in the Markdown file of its document at the commit 102f399:

{\footnotesize
\begin{longtable}{>{\raggedright\arraybackslash}p{\dimexpr 0.0700\linewidth-2\tabcolsep\relax}>{\raggedright\arraybackslash}p{\dimexpr 0.3300\linewidth-2\tabcolsep\relax}>{\raggedright\arraybackslash}p{\dimexpr 0.1400\linewidth-2\tabcolsep\relax}>{\raggedright\arraybackslash}p{\dimexpr 0.4600\linewidth-2\tabcolsep\relax}}
\toprule
\textbf{row} & \textbf{document} & \textbf{strings found} & \textbf{source} \\
\midrule\endfirsthead
\toprule
\textbf{row} & \textbf{document} & \textbf{strings found} & \textbf{source} \\
\midrule\endhead
\endfoot
\bottomrule\endlastfoot
2 & \texttt{README} & 1 of 1 & \texttt{git show 102f399:README.md} \\
3 & \texttt{DIRAC16COMPLEX\_\allowbreak{}PROVENANCE} & 2 of 2 & \texttt{git show 102f399:DIRAC16COMPLEX\_\allowbreak{}PROVENANCE.md} \\
4 & \texttt{GRAVITY\_\allowbreak{}PROVENANCE} & 2 of 2 & \texttt{git show 102f399:GRAVITY\_\allowbreak{}PROVENANCE.md} \\
5 & \texttt{STANDARD\_\allowbreak{}MODEL\_\allowbreak{}PROVENANCE} & 1 of 1 & \texttt{git show 102f399:STANDARD\_\allowbreak{}MODEL\_\allowbreak{}PROVENANCE.md} \\
6 & \texttt{COUPLED\_\allowbreak{}SOLUTIONS\_\allowbreak{}PROVENANCE} & 1 of 1 & \texttt{git show 102f399:COUPLED\_\allowbreak{}SOLUTIONS\_\allowbreak{}PROVENANCE.md} \\
7 & \texttt{COUPLED\_\allowbreak{}SOLUTIONS} & 3 of 3 & \texttt{git show 102f399:COUPLED\_\allowbreak{}SOLUTIONS.md} \\
8 & \texttt{COUPLED\_\allowbreak{}SOLUTIONS\_\allowbreak{}PROVENANCE} & 2 of 2 & \texttt{git show 102f399:COUPLED\_\allowbreak{}SOLUTIONS\_\allowbreak{}PROVENANCE.md} \\
9 & \texttt{DARK\_\allowbreak{}SECTOR} & 1 of 1 & \texttt{git show 102f399:DARK\_\allowbreak{}SECTOR.md} \\
10 & \texttt{DIRAC16COMPLEX} & 1 of 1 & \texttt{git show 102f399:DIRAC16COMPLEX.md} \\
11 & \texttt{DIRAC16COMPLEX\_\allowbreak{}PROVENANCE} & 1 of 1 & \texttt{git show 102f399:DIRAC16COMPLEX\_\allowbreak{}PROVENANCE.md} \\
14 & \texttt{KOHN\_\allowbreak{}SHAM\_\allowbreak{}GRAVITY\_\allowbreak{}PROVENANCE} & 2 of 2 & \texttt{git show 102f399:KOHN\_\allowbreak{}SHAM\_\allowbreak{}GRAVITY\_\allowbreak{}PROVENANCE.md} \\
15 & \texttt{README\_\allowbreak{}PROVENANCE} & 1 of 1 & \texttt{git show 102f399:README\_\allowbreak{}PROVENANCE.md} \\
16 & \texttt{STANDARD\_\allowbreak{}MODEL} & 2 of 2 & \texttt{git show 102f399:STANDARD\_\allowbreak{}MODEL.md} \\
17 & \texttt{STANDARD\_\allowbreak{}MODEL\_\allowbreak{}SYMBOLS} & 1 of 1 & \texttt{git show 102f399:STANDARD\_\allowbreak{}MODEL\_\allowbreak{}SYMBOLS.md} \\*
\end{longtable}}
\textbf{14 of 14} formulas found.

\Needspace{8\baselineskip}\subsection*{6.4 The checks against this build}\addcontentsline{toc}{subsection}{6.4 The checks against this build}
The exit codes and the counts (2.1) of the twelve programs, recounted from their outputs in this build of the documents, against the data file (exit code, True, False):

{\footnotesize
\begin{longtable}{>{\raggedright\arraybackslash}p{\dimexpr 0.4000\linewidth-2\tabcolsep\relax}>{\raggedright\arraybackslash}p{\dimexpr 0.2000\linewidth-2\tabcolsep\relax}>{\raggedright\arraybackslash}p{\dimexpr 0.2000\linewidth-2\tabcolsep\relax}>{\raggedright\arraybackslash}p{\dimexpr 0.2000\linewidth-2\tabcolsep\relax}}
\toprule
\textbf{program} & \textbf{data file} & \textbf{this build} & \textbf{agree} \\
\midrule\endfirsthead
\toprule
\textbf{program} & \textbf{data file} & \textbf{this build} & \textbf{agree} \\
\midrule\endhead
\endfoot
\bottomrule\endlastfoot
\texttt{tools/\allowbreak{}verify\_\allowbreak{}7p1\_\allowbreak{}nb.py} & 0, 54, 0 & 0, 54, 0 & yes \\
\texttt{tools/\allowbreak{}ref44.py} & 0, 10, 0 & 0, 10, 0 & yes \\
\texttt{tools/\allowbreak{}ref71.py} & 0, 25, 0 & 0, 25, 0 & yes \\
\texttt{tools/\allowbreak{}extra71.py} & 0, 23, 0 & 0, 23, 0 & yes \\
\texttt{tools/\allowbreak{}ref\_\allowbreak{}dirac16.py} & 0, 90, 0 & 0, 90, 0 & yes \\
\texttt{tools/\allowbreak{}ref\_\allowbreak{}gravity.py} & 0, 101, 0 & 0, 101, 0 & yes \\
\texttt{tools/\allowbreak{}ref\_\allowbreak{}ks\_\allowbreak{}gravity.py} & 0, 29, 0 & 0, 29, 0 & yes \\
\texttt{tools/\allowbreak{}ref\_\allowbreak{}sm.py} & 0, 33, 0 & 0, 33, 0 & yes \\
\texttt{tools/\allowbreak{}ref\_\allowbreak{}coupled.py} & 0, 50, 0 & 0, 50, 0 & yes \\
\texttt{tools/\allowbreak{}ref\_\allowbreak{}dark.py} & 0, 20, 0 & 0, 20, 0 & yes \\
\texttt{tools/\allowbreak{}ref\_\allowbreak{}torus\_\allowbreak{}brane.py} & 0, 38, 0 & 0, 38, 0 & yes \\
\texttt{tools/\allowbreak{}ref\_\allowbreak{}radion.py} & 0, 31, 0 & 0, 31, 0 & yes \\*
\end{longtable}}
\textbf{12 of 12} programs agree.

The notebook: the data file against the evaluated 7+1.nb as this build parses it, and against its evaluation record:

{\footnotesize
\begin{longtable}{>{\raggedright\arraybackslash}p{\dimexpr 0.5000\linewidth-2\tabcolsep\relax}>{\raggedright\arraybackslash}p{\dimexpr 0.2500\linewidth-2\tabcolsep\relax}>{\raggedright\arraybackslash}p{\dimexpr 0.2500\linewidth-2\tabcolsep\relax}}
\toprule
\textbf{quantity} & \textbf{data file} & \textbf{this build} \\
\midrule\endfirsthead
\toprule
\textbf{quantity} & \textbf{data file} & \textbf{this build} \\
\midrule\endhead
\endfoot
\bottomrule\endlastfoot
input cells & 189 & 189 \\
output cells & 554 & 554 \\
message cells & 0 & 0 \\
checks printed where they stand & 443 & 443 \\
all checks (with those of \texttt{frameChecks}) & 479 & 479 \\
checks that returned True & 479 & 479 \\
evaluation time (seconds) & 283.883 & 283.883 \\*
\end{longtable}}
\textbf{All seven agree.}

The comparison of the fresh clone of 102f399 with the commit (184 tracked files; 60 of 60 Markdown and LaTeX files identical; 30 of 30 PDFs with the same text and page count; 0 documents with layout problems; 0 other files changed) is the record of that check, which was made once, on 28 October 2026 at 13:35 PDT.

\Needspace{8\baselineskip}\subsection*{6.5 The formulas whose presence was checked}\addcontentsline{toc}{subsection}{6.5 The formulas whose presence was checked}
\Needspace{6\baselineskip}\subsubsection*{6.5.1 Row 2, README}
For the 256 increasing multi-indices $K=(A_1<\dots<A_r)$, $r=0,\dots,8$, of the eight real 16×16 Dirac matrices $\Gamma^A$:
\begin{equation*}
\Theta^K=i^{\sigma_r}\Gamma^{A_1}\cdots\Gamma^{A_r},\qquad \sigma_r=\frac{r(r-1)}2\bmod2,\qquad \Theta_K=\eta_{A_1A_1}\cdots\eta_{A_rA_r}\Theta^K .
\end{equation*}
The factor $i^{\sigma_r}$ makes each $\bar\Psi\Theta^K\Psi$ real; $\eta_{AA}$ are the diagonal entries of $\eta=\mathrm{diag}(+1,-1,\dots,-1)$; for $r=0$ the product is $1_{16}$.

\Needspace{6\baselineskip}\subsubsection*{6.5.2 Row 3, DIRAC16COMPLEX\_PROVENANCE}
For increasing index sets $K$, $K'$ of the anticommuting generators $\theta_k$ ($\theta_k\theta_l=-\theta_l\theta_k$), with $\theta_K$ the product of the generators of $K$ in increasing order:
\begin{align*}
\theta_K\theta_{K'}&=\mathrm{sign}(\pi)\,\theta_{\mathrm{sort}(K\cup K')}\qquad(\text{zero if }K\cap K'\neq\emptyset), \\
\frac{\partial^L(\theta_{k_1}\cdots\theta_{k_r})}{\partial\theta_{k_j}}&=(-1)^{j-1}\,\theta_{k_1}\cdots\widehat{\theta_{k_j}}\cdots\theta_{k_r} .
\end{align*}
$\pi$ is the permutation that sorts the sequence $K$ followed by $K'$, $\mathrm{sort}(K\cup K')$ the sorted union, $\partial^L$ the derivative from the left, and the hat marks the factor that is left out.

\Needspace{6\baselineskip}\subsubsection*{6.5.3 Row 4, GRAVITY\_PROVENANCE}
As the author’s notebook Einstein-Lovelock-7+1-Nash.nb writes them, after a component metric $g$ has been assigned (Wolfram Language with the packages xAct, xCoba and xTras):

\begin{Verbatim}[fontsize=\small,frame=single,framesep=2mm]
CD = LC[g]; RiemannCD = Riemann[CD]; RicciCD = Ricci[CD];
RicciScalarCD = RicciScalar[CD]; EinsteinCD = Einstein[CD];
\end{Verbatim}
After \texttt{CD = LC[g]} the abstract covariant derivative \texttt{CD} is the Levi-Civita connection of the component metric $g$, and \texttt{RiemannCD}, \texttt{RicciCD}, \texttt{RicciScalarCD}, \texttt{EinsteinCD} are its Riemann tensor, Ricci tensor, scalar curvature and Einstein tensor.

\Needspace{6\baselineskip}\subsubsection*{6.5.4 Row 5, STANDARD\_MODEL\_PROVENANCE}
Diagonalizing the two-flavour Hamiltonian in matter gives the mixing angle in matter $\theta_m$:
\begin{equation*}
\frac{\Delta m^2}{4E}\begin{pmatrix}-\cos2\theta&\sin2\theta\\\sin2\theta&\cos2\theta\end{pmatrix}+\mathrm{diag}(V,0)\quad\Longrightarrow\quad \tan2\theta_m=\frac{\sin2\theta}{\cos2\theta-2EV/\Delta m^2} .
\end{equation*}
$\Delta m^2$ is the difference of the squared masses, $E$ the neutrino energy, $\theta$ the mixing angle in vacuum and $V$ the potential of the electron neutrino in matter.

\Needspace{6\baselineskip}\subsubsection*{6.5.5 Row 6, COUPLED\_SOLUTIONS\_PROVENANCE}
The two-component amplitude $c(\tau)$ of a mode of dirac16complex in the Einstein–Kasner solution of the earliest regime obeys
\begin{equation*}
i\,\frac{dc}{d\tau}=\omega\,(\cos\theta\,\sigma_1+\sin\theta\,\sigma_2)\,c,\qquad \omega=K\sqrt{\tau^{-2p_1}+\tau^{-2q_1}},\qquad \tan\theta=\tau^{p_1-q_1} .
\end{equation*}
$\tau$ is the time since the singularity, $K$ the scaling variable of the mode, $p_1=(1+2\sqrt2)/7$ and $q_1=u_1p_1$ with $u_1=-1+1/\sqrt2$ the Kasner exponents of ordinary and auxiliary space, and $\sigma_1$, $\sigma_2$ Pauli matrices.

\Needspace{6\baselineskip}\subsubsection*{6.5.6 Row 7, COUPLED\_SOLUTIONS}
For the metric $ds^2=n^2dx_0^2-a^2(dx_1^2+dx_2^2+dx_3^2)-b^2(dx_4^2+\dots+dx_7^2)$, with $\mathrm{EL}_q$ the Euler–Lagrange expression for the variable $q$, evaluated at $n=1$:
\begin{align*}
S&=-\frac1{2\kappa}\int n\,a^3b^4\sum_k\alpha_kL_k\,dx_0-\int nW\,dx_0, \\
\mathrm{EL}_n(na^3b^4L_k)&=-2a^3b^4E^{(k)0}{}_0,\qquad \mathrm{EL}_a(na^3b^4L_k)=-6a^2b^4E^{(k)1}{}_1, \\
\mathrm{EL}_b(na^3b^4L_k)&=-8a^3b^3E^{(k)4}{}_4 .
\end{align*}
$\kappa=8\pi G_8$, $\alpha_k$ are the Lovelock couplings ($\alpha_0=2\Lambda$, $\alpha_1=1$), $L_k$ the Lovelock densities, $E^{(k)A}{}_B$ the Lovelock tensors and $W(a,b)$ the comoving energy of the matter.

\Needspace{6\baselineskip}\subsubsection*{6.5.7 Row 8, COUPLED\_SOLUTIONS\_PROVENANCE}
For one massless particle with comoving momenta $k_a$ along ordinary space and $k_b$ along the auxiliary space:
\begin{equation*}
W=\omega=\sqrt{k_a^2/a^2+k_b^2/b^2},\qquad \partial_aW=-\frac{k_a^2}{a^3\omega},\qquad \partial_bW=-\frac{k_b^2}{b^3\omega} .
\end{equation*}
\begin{align*}
S&=-\frac1{2\kappa}\int na^3b^4\sum_k\alpha_kL_k\,dx_0-\int nW\,dx_0, \\
\mathrm{EL}_n(na^3b^4L_k)&=-2a^3b^4E^{(k)0}{}_0,\qquad \mathrm{EL}_a(na^3b^4L_k)=-6a^2b^4E^{(k)1}{}_1, \\
\mathrm{EL}_b(na^3b^4L_k)&=-8a^3b^3E^{(k)4}{}_4 .
\end{align*}
The second display is the minisuperspace action with the identities of its variation, the same equations as those of the previous paragraph, except that the provenance file writes $na^3b^4$ without the thin space of $n\,a^3b^4$; $\omega$ is the energy of the particle, and the derivatives $\partial_aW$, $\partial_bW$ enter the variations with respect to $a$ and $b$.

\Needspace{6\baselineskip}\subsubsection*{6.5.8 Row 9, DARK\_SECTOR}
A level of Kaluza–Klein mass $\mu_n$ of dirac16complex (mass $m$) has, in the Einstein frame, the mass $M_E(\beta)$, with $M_n=M_E(0)$, and hence the coupling $g_n$ to the canonical radion $\phi=\sqrt{12}M_{\rm Pl}\beta$:
\begin{equation*}
M_E(\beta)=e^{-2\beta}\sqrt{m^2+\mu_n^2e^{-2\beta}},\qquad g_n=\Big|\frac{dM_E}{d\phi}\Big|_{\beta=0}=\frac{2M_n+\mu_n^2/M_n}{\sqrt{12}\,M_{\rm Pl}} .
\end{equation*}
$\beta=\ln(b/b_0)$ is the radion and $M_{\rm Pl}$ the reduced Planck mass.

\Needspace{6\baselineskip}\subsubsection*{6.5.9 Row 10, DIRAC16COMPLEX}
7+1.nb sums the mode algebra exactly, with symbolic $\vec x,\vec y$ and $m=1$, for the symmetric set of five momenta
\begin{equation*}
\vec p\in\big\{\pm(2,2,4,0,0,0,0),\ \pm(0,0,0,2,2,0,0),\ \vec 0\big\}\qquad(E=5,5,3,3,1).
\end{equation*}
$E=\sqrt{m^2+|\vec p|^2}$: $\sqrt{1+24}=5$, $\sqrt{1+8}=3$, $\sqrt1=1$.

\Needspace{6\baselineskip}\subsubsection*{6.5.10 Row 11, DIRAC16COMPLEX\_PROVENANCE}
With the real spinors $u_s(\vec p)$ of positive and $v_s(-\vec p)$ of negative energy ($s=1,\dots,8$, $u_s^{\mathsf T}u_s=2E$) and the real symmetric plane-wave Dirac Hamiltonian $h(\vec p)$:
\begin{equation*}
U=\frac1{\sqrt{2E}}\big(u_1(\vec p),\dots,u_8(\vec p),v_1(-\vec p),\dots,v_8(-\vec p)\big)\qquad\text{satisfies}\qquad U^{\mathsf T}h(\vec p)\,U=\mathrm{diag}(E1_8,-E1_8).
\end{equation*}
$U$ is real and orthogonal, $U^{\mathsf T}U=1_{16}$; $h(\vec p)=p^i\alpha_i+m\Gamma^0$ with $\alpha_i=\Gamma^0\Gamma^i$ ($i=1,\dots,7$, summed).

\Needspace{6\baselineskip}\subsubsection*{6.5.11 Row 14, KOHN\_SHAM\_GRAVITY\_PROVENANCE}
The 256 matrices $\Theta_K$ are complete (the first display), and this gives the Fierz identity (the second display) for commuting spinors $a,b,c,d$:
\begin{equation*}
\sum_K(\Theta_K)_{cd}(\Theta^K)_{ab}=16\,\delta_{ad}\delta_{bc},
\end{equation*}
\begin{equation*}
\sum_{K\in r}(a\Theta_Kb)(c\Theta^Kd)=\sum_JC_{rJ}\,(a\Theta_Jd)(c\Theta^Jb),\qquad C_{rJ}=\frac1{256}\sum_{K\in r}\mathrm{tr}\big(\Theta_K\Theta_J\Theta^K\Theta^J\big).
\end{equation*}
Here $(aMb)=\sum_{i,j}a_iM_{ij}b_j$ for a 16×16 matrix $M$; $K\in r$ runs over the multi-indices of rank $r$ and $J$ over all 256; as subscripts in the first display, $a,b,c,d$ are spinor indices $1,\dots,16$, and $\delta$ is the Kronecker delta.

\Needspace{6\baselineskip}\subsubsection*{6.5.12 Row 15, README\_PROVENANCE}
With $\Psi$ the field dirac16complex:
\begin{gather*}
\Gamma_A=\Gamma^{(1,7)}_A,\qquad \Gamma^{A}=\eta_{(1,7)}^{AA}\Gamma_A,\qquad \Gamma^0=\Gamma_0,\qquad \chi=\chi^{(1,7)}=i\Gamma^{(1,7)}_8,\\ \bar\Psi=\Psi^\dagger\Gamma^0,\qquad \vec x=(x_1,\dots,x_7),\qquad \vec p=(p^1,\dots,p^7).
\end{gather*}
$\Gamma^{(1,7)}_A$ are the eight real 16×16 Dirac matrices of 7+1.nb, $\eta_{(1,7)}=\mathrm{diag}(+1,-1,\dots,-1)$, and $\Gamma^{(1,7)}_8=\Gamma^{(1,7)}_0\cdots\Gamma^{(1,7)}_7$.

\Needspace{6\baselineskip}\subsubsection*{6.5.13 Row 16, STANDARD\_MODEL}
The reworded sentence states:
\begin{equation*}
v=(\sqrt2G_F)^{-1/2}=246.2196402\ \text{GeV},\qquad m_W=\frac{gv}2,\qquad m_Z=\frac{\sqrt{g^2+g'^2}\,v}2,\qquad m_H^2=2\lambda v^2 .
\end{equation*}
$G_F$ is the Fermi constant, $g$ and $g'$ the gauge couplings of $SU(2)_L$ and $U(1)_Y$, $\lambda$ the quartic coupling of the Higgs field; the document writes the three masses inline, as $m_W=gv/2$ and $m_Z=\sqrt{g^2+g'^2}\,v/2$.

\Needspace{6\baselineskip}\subsubsection*{6.5.14 Row 17, STANDARD\_MODEL\_SYMBOLS}
The reworded convention:
\begin{gather*}
1_n=(\delta_{ij})_{i,j=1}^n,\qquad (M^\dagger)_{ij}=\overline{M_{ji}},\qquad (M^*)_{ij}=\overline{M_{ij}},\\ (A\otimes B)_{(i-1)p+k,\,(j-1)q+l}=A_{ij}B_{kl}\ \ (B\ \text{of size}\ p\times q).
\end{gather*}
$1_n$ is the $n\times n$ identity, $M^\dagger$ the conjugate transpose, $M^*$ the complex conjugate (the bar is complex conjugation of a number) and $A\otimes B$ the Kronecker product, in Mathematica \texttt{KroneckerProduct[A, B]}; the document states these in words, the display writes them out.

\Needspace{12\baselineskip}\section*{7. Files}\addcontentsline{toc}{section}{7. Files}
\textbf{SHA-256 of the files} (at the time this record was generated):

\begin{Verbatim}[fontsize=\small,frame=single,framesep=2mm]
7plus1_status_2026-10-28.xlsx  d1624550cb0eefb3ead87a56f0055daa96bbd37af37b98434a60bf31f282f315
tools/status_2026-10-28.json   653435c67807abbcab805a369fcf75a307dc27185ad60d3977316d898071f715
tools/make_status_xlsx.py      53e063161fcb9acc88749d491c26398e499f22da46e64792a2c828caecd19a47
tools/doc_status_common.py     0e81361305921d951b7d614eb49eb226bd225acca9f30820cf80d0d5fa71e222
tools/doc_status.py            a8f940cb0fea5bc88381e77bba1e9c6d7ca8992dd856688dd01620cd39cfc2d5
tools/doc_status_symbols.py    c9720ccf538b5c657bb54ceb53e5285b005a9fdd113dd485018622033a6ec1b7
tools/doc_status_provenance.py 059bbcda81f46a9bb9a968e072ca6867c3e6aaa53940ad1475142451594eec5b
tools/make_docs.py             88998b5952ce99e0e04f65caf9959eefde0658e4dcfc73d936818e2e49c3b400
tools/docgen.py                747d5abe5347077802cd302a08ad293e5b458630f73ce4f947c0686268609ca2
tools/doc_common.py            51995b299ba2aa0947fcbfd593486cb9f2cd724d15f7a43035f7a47f8ade454c
\end{Verbatim}
The first version of the spreadsheet, which is not in the repository, had the SHA-256 \texttt{179f155a\allowbreak{}c681ae0f\allowbreak{}3d7d283d\allowbreak{}c1e427e9\allowbreak{}dd4badbf\allowbreak{}2b3f2198\allowbreak{}090059f2\allowbreak{}a38d0b83}.

\end{document}
